\scanpage{165}
\deleted{\unsure{Soit maintenant}} \inserted{Plus gén\textsuperscript{t} soit}
$M$ un motif effectif, \unsure{hom} de degré $i$.
Choisissons un hom \uline{injectif} $M\longrightarrow\uline h^i(X)$
et considérons la forme induite par une $\varphi\in\Phi_X$, soit $\varphi_0$.
Choisissant de $=$ $M\longrightarrow\uline h^i(Y)$, et $\psi\in\Phi_Y$,
d'où $\psi_0$. Je dis que $\varphi_0$ et $\psi_0$ sont
\inserted{\uline{définies}} et \uline{équivalentes}.
\deleted{ce qui signifie que $\operatorname{Tr}$}
La première p\textsuperscript{té} résulte des
\deleted{\unsure{propriétés}} \deleted{\unsure{le fait que} $\varphi$ est définie}
\deleted{\unsure{les}} \deleted{positive}\ldots Riemanniennes\ldots
La deuxième \deleted{\unsure{reste}}

\diagram{\begin{tikzcd}[column sep=small,row sep=large]
\substack{\uline h^i(X)\\ \varphi}
  \arrow[dr,"\lambda"]
  \arrow[dr,bend right=35,"{\deleted{\unsure{\varphi|\psi}}\lambda'}"']
& & & \substack{\uline h^{\unsure{j|i}}(Y)\\ \psi}
  \arrow[dl,"\mu"'] \arrow[dl,bend left=35,"\mu'"]\\
& \substack{M\\ \varphi_0}\arrow[r,"u"] & \substack{M\\ \psi_0} &
\end{tikzcd}}

exprime que la forme $\operatorname{Tr}(u'u)$ est définie,
\inserted{$u\in\operatorname{End}(M)$}
où $u'$ est définie par $\varphi_0$ et $\psi_0$. Or
\deleted{\unsure{car}} pour cel, il suffit de
\deleted{\unsure{prouver}} \inserted{noter} que la forme
\[
\operatorname{Tr}(\mu u\lambda')'(\mu u\lambda')
=\operatorname{Tr}\lambda u'\mu'\mu u\lambda'
=\operatorname{Tr}\lambda'\lambda u'\mu'\mu u
\]
l'est, or $\mu'\mu=\mathrm{id}$,
\inserted{$\lambda'\lambda=\mathrm{id}$}
donc cette forme est $\operatorname{Tr}u'u$, OK !
On pourra donc, \uline{sur un motif effectif}, définir une classe
d'équivalence de formes riemanniennes, \unsure{soit}
$\unsure{\Phi_M|\Psi_M}$ (\uline{\unsure{ens.} saturé !}). Les $\Psi_M$ satisfont
\mspage{164}
